%======================================================================
\section{Constraining the Shadow: Regularisation}\label{sec:regularisation}
%======================================================================

Until now, every result has been a consequence of orthogonal projection.
The hat matrix $\bH$ is idempotent and symmetric; the fitted values
$\yhat$ are the closest point in $\Col(\bX)$ to $\by$; the residual
$\be$ meets the column space at a right angle.
Regularisation breaks this pattern---Ridge and LASSO are \emph{not}
projections.
They deliberately sacrifice the ``closest point'' property to gain
stability.
This section explains how and why.

\begin{warning}
\textbf{Where the projection framework reaches its limits.}
The Ridge hat matrix
$\bH_{\text{ridge}} = \bX(\bX^\top\bX+\lambda\bI)^{-1}\bX^\top$
is \emph{not} idempotent
($\bH_{\text{ridge}}^2 \neq \bH_{\text{ridge}}$) and is \emph{not}
an orthogonal projection onto any subspace.
It does not ``drop a perpendicular''---it \emph{pulls the shadow
inward}, toward the origin, away from the foot of the perpendicular.

The LASSO solution is not even linear in $\by$.

The geometric language of shadows and right angles served us well for
OLS\@.  For regularisation, we need a different picture: that of
\textbf{constrained optimisation}, where the solution is determined
not by orthogonality but by the tangency between an objective function's
level sets and a constraint region.
\end{warning}

\medskip
\noindent\textbf{The telescope analogy.}
Think of the OLS estimate as looking through a telescope on a windy night.
The telescope points in the right direction on average (unbiasedness),
but the wind makes it shake violently (high variance from
multicollinearity or small sample size).
Each snapshot is centred on the truth, but individual frames are blurry
and useless.

Regularisation is like tightening the telescope's mount.
A tighter mount introduces a small systematic error---the telescope
now points slightly away from the true target (bias).
But the shaking is dramatically reduced (lower variance), so each
individual frame is much sharper and more useful.

Ridge tightens the mount uniformly in all directions.
LASSO tightens some axes completely (locking them to zero) while
leaving others free---it is like a mount that can only swivel along
certain axes, automatically selecting which directions matter.

The bias-variance trade-off is the trade-off between pointing accuracy
and shake reduction.
The optimal penalty $\lambda$ is the tightness that minimises the
total blur (mean squared error $= \text{bias}^2 + \text{variance}$).


%----------------------------------------------------------------------
\subsection{Ridge regression}
%----------------------------------------------------------------------

\begin{definition}
The \textbf{Ridge estimator} with penalty $\lambda\ge 0$ is
\[
  \bhat_{\text{ridge}}
  = \argmin_{\bm{\beta}}\left\{
      \norm{\by-\bX\bm{\beta}}^2 + \lambda\norm{\bm{\beta}}^2
    \right\}
  = (\bX^\top\bX + \lambda\bI)^{-1}\bX^\top\by.
\]
\end{definition}

\noindent\textbf{Eigenvalue shrinkage---the key mechanism.}
Ridge adds $\lambda\bI$ to $\bX^\top\bX$, which shifts every eigenvalue
$\lambda_k\mapsto\lambda_k+\lambda$.
Write the eigendecomposition $\bX^\top\bX = \bQ\bD\bQ^\top$, where
$\bD = \diag(\lambda_1,\dots,\lambda_p)$.
Then
\begin{equation}\label{eq:ridge-eigen}
  \bhat_{\text{ridge}}
  = \bQ\,\diag\!\left(\frac{\lambda_k}{\lambda_k+\lambda}\right)\bQ^\top
    \bhat_{\text{OLS}}.
\end{equation}
Each eigenvector component of $\bhat_{\text{OLS}}$ is multiplied by the
shrinkage factor $\lambda_k/(\lambda_k+\lambda)$, which lies in $[0,1]$.
\begin{itemize}[nosep]
  \item Directions with \emph{large} $\lambda_k$ (stable, well-determined
        by the data) are barely affected: the factor is close to $1$.
  \item Directions with \emph{small} $\lambda_k$ (unstable, dominated by
        noise) are aggressively shrunk toward zero: the factor is close
        to $0$.
\end{itemize}

\begin{keyinsight}
Ridge performs \emph{selective geometric surgery}: it identifies the
unstable directions of the column space via the eigenvalue spectrum
and shrinks those directions most aggressively, while leaving the
stable directions nearly intact.
This is \emph{anisotropic} shrinkage---it does not pull equally in
all directions, but preferentially collapses the dimensions where OLS
is least trustworthy.
\end{keyinsight}


%--- Example 6: Ridge shrinkage direction by direction -----------------

\begin{example}[Ridge Shrinkage Direction by Direction]\label{ex:ridge}
We work with a $2\times 2$ cross-product matrix $\bX^\top\bX$ whose
eigenvalues and eigenvectors are known, to see exactly how Ridge shrinks
each direction.

\medskip\noindent\textbf{Setup.}
Suppose (after centring) $\bX^\top\bX$ has eigendecomposition
$\bX^\top\bX = \bQ\bD\bQ^\top$ with eigenvalues $9$ and $1$:
\[
\bQ = \frac{1}{\sqrt{2}}\begin{pmatrix}1&1\\1&-1\end{pmatrix},\quad
\bD = \begin{pmatrix}9&0\\0&1\end{pmatrix}.
\]
This gives
\[
\bX^\top\bX
= \frac{1}{2}\begin{pmatrix}1&1\\1&-1\end{pmatrix}
\begin{pmatrix}9&0\\0&1\end{pmatrix}
\begin{pmatrix}1&1\\1&-1\end{pmatrix}
= \begin{pmatrix}5&4\\4&5\end{pmatrix}.
\]
The ``wide'' eigendirection is
$\bm{v}_1 = (1,1)^\top/\sqrt{2}$ (eigenvalue $9$, lots of data spread)
and the ``narrow'' eigendirection is
$\bm{v}_2 = (1,-1)^\top/\sqrt{2}$ (eigenvalue $1$, little spread).

Choose OLS solution $\bhat_{\text{OLS}} = (3,\;1)^\top$.
Decompose in the eigenbasis using unnormalised directions
$\bm{u}_1 = (1,1)^\top$, $\bm{u}_2 = (1,-1)^\top$:
\[
(3,1) = \alpha_1(1,1) + \alpha_2(1,-1)
\implies \alpha_1 = 2,\;\alpha_2 = 1.
\]
So OLS has component $2$ along the stable direction and $1$ along
the unstable direction.

\medskip\noindent\textbf{Ridge coefficients.}
In the eigenbasis, each component is simply multiplied by a shrinkage
factor:
\[
\alpha_k^{\text{Ridge}}
= \frac{\lambda_k}{\lambda_k + \lambda}\;\alpha_k^{\text{OLS}}.
\]
No matrix inversion is needed once you work in the eigenbasis.

\begin{center}
\renewcommand{\arraystretch}{1.4}
\begin{tabular}{c|cc|cc|cc|c}
\toprule
& \multicolumn{2}{c|}{Shrinkage factor}
& \multicolumn{2}{c|}{Eigencomponents}
& \multicolumn{2}{c|}{Coefficients} & \\
$\lambda$
& $\frac{9}{9+\lambda}$ & $\frac{1}{1+\lambda}$
& $\alpha_1^{\text{R}}$ & $\alpha_2^{\text{R}}$
& $\hat\beta_1^{\text{R}}$ & $\hat\beta_2^{\text{R}}$
& $\norm{\bhat_{\text{R}}}$ \\
\midrule
$0$   & $1$ & $1$
      & $2$ & $1$
      & $3$ & $1$
      & $\sqrt{10}\approx 3.16$ \\[2pt]
$1$   & $\frac{9}{10}$ & $\frac{1}{2}$
      & $\frac{9}{5}$ & $\frac{1}{2}$
      & $\frac{23}{10}$ & $\frac{13}{10}$
      & $2.64$ \\[2pt]
$10$  & $\frac{9}{19}$ & $\frac{1}{11}$
      & $\frac{18}{19}$ & $\frac{1}{11}$
      & $1.04$ & $0.86$
      & $1.35$ \\[2pt]
$100$ & $\frac{9}{109}$ & $\frac{1}{101}$
      & $\frac{18}{109}$ & $\frac{1}{101}$
      & $0.18$ & $0.16$
      & $0.24$ \\
\bottomrule
\end{tabular}
\end{center}

\medskip\noindent\textbf{The key pattern.}
At $\lambda=1$, the stable direction (eigenvalue $9$) retains $90\%$ of
its OLS component, while the unstable direction (eigenvalue $1$) is
already halved.
At $\lambda=10$, the unstable direction is crushed to $9\%$ of its OLS
value, while the stable direction still retains $47\%$.
As $\lambda$ increases, the coefficient vector traces a \emph{curved}
path from $(3,1)$ toward the origin---not a straight line, because
the unstable eigencomponent collapses first.
At $\lambda=100$, both coefficients are nearly equal
($0.18$ vs.\ $0.16$) because only the stable ``sum'' direction
retains any signal.
\end{example}


%--- Stop and Think: Ridge limits --------------------------------------

\begin{stopandthink}
\textbf{Ridge limits.}

Recall that
$\bhat_{\text{ridge}} = (\bX^\top\bX + \lambda\bI)^{-1}\bX^\top\by$,
and in the eigendecomposition each component is shrunk by
$\lambda_k / (\lambda_k + \lambda)$.
\begin{enumerate}[nosep]
  \item What happens as $\lambda \to 0$?
        What does each shrinkage factor approach?
  \item What happens as $\lambda \to \infty$?
        What does each shrinkage factor approach?
        What does $\bhat_{\text{ridge}}$ approach?
  \item Is there a value of $\lambda$ for which Ridge is
        equivalent to ignoring the data entirely?
\end{enumerate}
\emph{Answers:}
(1)~As $\lambda \to 0$, each factor
$\lambda_k/(\lambda_k + \lambda) \to 1$,
so $\bhat_{\text{ridge}} \to \bhat_{\text{OLS}}$.
Ridge with no penalty is OLS\@.

(2)~As $\lambda \to \infty$, each factor
$\lambda_k/(\lambda_k + \lambda) \to 0$,
so $\bhat_{\text{ridge}} \to \bm{0}$.
Infinite penalty shrinks all coefficients to zero---the model predicts
$\hat{y}_i = 0$ for everyone (or $\hat{y}_i = \bar{y}$ if the intercept
is not penalised).

(3)~Yes: $\lambda = \infty$ means the model ignores the data
completely and defaults to the prior that all coefficients are zero.
Ridge traces a continuous path from ``trust the data fully''
($\lambda=0$) to ``trust the prior fully'' ($\lambda=\infty$).
\end{stopandthink}


\medskip
\noindent\textbf{Geometric interpretation.}
The OLS solution must lie in the unconstrained parameter space.
The Ridge solution must lie inside a sphere
$\{\bm{\beta}:\norm{\bm{\beta}}\le t\}$ for some $t$ determined
by $\lambda$.
If the OLS solution is already inside the sphere, Ridge and OLS agree.
Otherwise, the Ridge solution is the point on the sphere's surface
that minimises the residual sum of squares---where the RSS
contour ellipse first touches the sphere.
See Figure~\ref{fig:ridge-lasso} for this picture in two dimensions:
the solid circle is the $L^2$ constraint region, and the Ridge
solution is the tangency point between the RSS ellipse and the sphere.


%----------------------------------------------------------------------
\subsection{LASSO regression}
%----------------------------------------------------------------------

\begin{definition}
The \textbf{LASSO estimator} with penalty $\lambda\ge 0$ is
\[
  \bhat_{\text{lasso}}
  = \argmin_{\bm{\beta}}\left\{
      \norm{\by-\bX\bm{\beta}}^2 + \lambda\norm{\bm{\beta}}_1
    \right\},
\]
where $\norm{\bm{\beta}}_1 = \sum_j|\beta_j|$.
\end{definition}

\noindent\textbf{Corner geometry and sparsity.}
The $L^1$ constraint region
$\{(\beta_1,\dots,\beta_p) : |\beta_1|+\cdots+|\beta_p| \le t\}$
is a diamond (cross-polytope) instead of a sphere.
Its surface has \emph{corners}---points where one or more coordinates
are exactly zero.
The RSS contour ellipses are more likely to first touch the diamond
at a corner, producing \textbf{exact zeros} in the coefficient vector
(see Figure~\ref{fig:ridge-lasso}, dashed diamond).

Why?
A corner is a ``pointy'' protrusion: an expanding ellipse approaching
from the direction of $\bhat_{\text{OLS}}$ will generically hit a
corner before it hits a smooth edge.
In contrast, the $L^2$ sphere has no corners, so the tangency point
generically has \emph{all} coordinates nonzero.

\medskip
\noindent\textbf{Sparsity as a corner solution.}
LASSO performs automatic variable selection: as $\lambda$ increases,
coefficients are driven to zero one by one.
Each zero corresponds to the solution hitting a corner of the $L^1$
diamond.
This is not an algebraic trick---it is a consequence of the geometry
of touching a polytope with an ellipsoid.

\begin{stopandthink}
\textbf{Why corners attract LASSO solutions.}

Grab a piece of paper and try this:
\begin{enumerate}[nosep]
  \item In $\R^2$, draw the $L^1$ diamond
        $\{(\beta_1, \beta_2) : |\beta_1| + |\beta_2| \le 1\}$.
        It has corners at $(1,0)$, $(-1,0)$, $(0,1)$, $(0,-1)$.
  \item Draw a family of concentric ellipses centred at some point
        outside the diamond (this represents the RSS contours centred
        at $\bhat_{\text{OLS}}$).
  \item Starting from the smallest ellipse, gradually expand it until
        it first touches the diamond. Where does it touch?
  \item Now replace the diamond with the $L^2$ circle
        $\{(\beta_1, \beta_2) : \beta_1^2 + \beta_2^2 \le 1\}$
        and repeat. Where does the first contact occur now?
\end{enumerate}
\emph{Key observation:}
The diamond's corners are ``pointy''---an expanding ellipse is far more
likely to first touch a corner (where one coordinate is exactly zero)
than a smooth edge.
The circle has no corners, so the first contact is almost always at a
point where \emph{both} coordinates are nonzero.
This is why LASSO produces sparse solutions (exact zeros) and Ridge
does not.
The geometry of the constraint set, not the algebra, is what drives
sparsity.
\end{stopandthink}


%----------------------------------------------------------------------
\subsection{Ridge vs.\ LASSO: a comparison}
%----------------------------------------------------------------------

Figure~\ref{fig:ridge-lasso} illustrates both constraint regions
simultaneously.
The OLS estimate $\bhat_{\text{OLS}}$ sits at the centre of the
elliptical RSS contours.
Ridge constrains the solution to the $L^2$ sphere (solid circle);
LASSO constrains it to the $L^1$ diamond (dashed).
The regularised solutions are the tangency points.

\begin{center}
\begin{tabular}{lcc}
\toprule
& \textbf{Ridge ($L^2$)} & \textbf{LASSO ($L^1$)} \\
\midrule
Constraint shape & Sphere (smooth) & Diamond (corners) \\
Shrinkage type & Anisotropic via eigenvalues & Along coordinate axes \\
Coefficient path & Smooth shrinkage toward 0 & Coefficients hit 0 exactly \\
Variable selection & No (all coefficients remain) & Yes (sparsity) \\
Best when & Many small effects (dense) & Few large effects (sparse) \\
Closed-form solution & Yes & No (requires iterative solver) \\
\bottomrule
\end{tabular}
\end{center}


%----------------------------------------------------------------------
\subsection{Reconnecting to the geometric framework}
%----------------------------------------------------------------------

Neither Ridge nor LASSO is an orthogonal projection.
Both sacrifice the ``closest point'' property---accepting bias---in
exchange for a (potentially large) reduction in variance.
The telescope analogy captures the trade-off:
you give up perfect aim to eliminate the shake.

Yet the connection to the projection framework runs deeper than it
first appears.
Recall from Section~\ref{sec:multicollinearity} that
multicollinearity makes the column space geometrically thin:
the eigenvalue spectrum of $\bX^\top\bX$ has a few large values
and some dangerously small ones.
OLS faithfully projects onto this flattened space, and in doing so
amplifies noise along the thin directions.
Ridge responds by \emph{reshaping} the eigenvalue landscape:
it replaces $\lambda_k$ with $\lambda_k + \lambda$, inflating every
eigenvalue and rounding out the column space's geometry.
The thinner a direction was, the more it gets inflated in relative
terms---precisely the directions that needed stabilising.

Ridge is the projection framework's answer to its own weakness.
When the geometry is unstable, you reshape the geometry.
